\begin{table}[H]
\centering
\caption{Νομοθετικές πράξεις που θεμελίωσαν το \ENG{ADL}}
\label{tab:adl_legislation}
\small
\setlength{\tabcolsep}{5pt}
\renewcommand{\arraystretch}{1.25}
\begin{chaptertablebox}
\begin{tabular}{|p{5cm}|p{8.5cm}|}
\hline
\textbf{Πράξη} & \textbf{Συμβολή} \\
\hline
\hline
\textbf{\ENG{Public} \ENG{Law} 105-261} (\ENG{Strom} \ENG{Thurmond} \ENG{National} \ENG{Defense} \ENG{Authorization} \ENG{Act}, 1999) & Παρείχε νομική βάση και χρηματοδότηση για την ανάπτυξη κοινών τεχνολογικών προτύπων εκπαίδευσης στο \ENG{DoD}~\cite{congress1999ndaa}. \\
\hline
\textbf{\ENG{Executive} \ENG{Order} 13111} (1999) & Επέκτεινε την εφαρμογή του \ENG{distributed} \ENG{learning} σε ολόκληρη την ομοσπονδιακή κυβέρνηση, καθιστώντας το \ENG{ADL} πανομοσπονδιακό πρόγραμμα~\cite{whitehouse1999eo13111}. \\
\hline
\textbf{\ENG{Department} \ENG{of} \ENG{Defense} \ENG{Strategic} \ENG{Plan} \ENG{for} \ENG{ADL}} (1999) & Όρισε τους στρατηγικούς στόχους: διαλειτουργικότητα, αποθετήρια επαναχρησιμοποιήσιμου περιεχομένου και ενσωμάτωση παιδαγωγικών προσεγγίσεων~\cite{dod1999strategic}. \\
\hline
\end{tabular}
\end{chaptertablebox}
\end{table}
